\DocumentMetadata{
  pdfversion=2.0,
  pdfstandard=ua-2,
  lang=en-US,
  tagging=on,
  tagging-setup={math/setup=mathml-SE,tabsorder=structure}
}
\documentclass[11pt,letterpaper]{article}
\usepackage[margin=0.68in,includefoot]{geometry}
\usepackage{newcomputermodern}
\usepackage{amsmath,amscd,mathtools}
\usepackage{tikz,adjustbox}
\providecommand{\square}{\mdlgwhtsquare}
\usepackage[hidelinks]{hyperref}
\hypersetup{
  pdftitle={Algebra First-Year Exam, May 2016, Part 1},
  pdfauthor={University of Florida Department of Mathematics},
  pdfsubject={Graduate mathematics examination},
  pdfdisplaydoctitle=true,
  bookmarksopen=true
}
\setcounter{secnumdepth}{0}
\setlength{\parindent}{0pt}
\setlength{\parskip}{0.28em}
\setlength{\emergencystretch}{3em}
\setlength{\leftmargini}{2.15em}
\setlength{\leftmarginii}{2.65em}
\setlength{\leftmarginiii}{3em}
\renewcommand{\labelenumi}{\textbf{\theenumi.}}
\pagestyle{empty}
\raggedbottom
\allowdisplaybreaks
\newenvironment{examproblems}[1]{%
  \begin{enumerate}%
  \setcounter{enumi}{#1}%
  \setlength{\topsep}{0.35em}%
  \setlength{\partopsep}{0pt}%
  \setlength{\itemsep}{0.48em}%
  \setlength{\parsep}{0.18em}%
}{\end{enumerate}}
\newenvironment{examsubparts}{%
  \begin{enumerate}%
  \setlength{\topsep}{0.18em}%
  \setlength{\partopsep}{0pt}%
  \setlength{\itemsep}{0.16em}%
  \setlength{\parsep}{0.1em}%
}{\end{enumerate}}
\newenvironment{examinstructions}{%
  \par\begingroup\itshape
}{%
  \par\endgroup\vspace{0.18em}
}
\makeatletter
\renewcommand\section{\@startsection{section}{1}{0pt}%
  {-1.25ex plus -.25ex minus -.1ex}%
  {0.55ex plus .1ex}%
  {\normalfont\large\bfseries}}
\renewcommand{\@maketitle}{%
  \newpage\null\vskip -1em%
  {\footnotesize\bfseries
    \href{https://www.ufl.edu/}{University of Florida}, \href{https://math.ufl.edu/}{Department of Mathematics}\par}%
  \vskip -0.5em%
  \tagstructbegin{tag=Title}%
    {\LARGE\bfseries\href{https://gma.math.ufl.edu/exams/algebra-first-year/}{Algebra First-Year Exam}, May 2016, Part 1\par}%
  \tagstructend%
  \vskip 0.25em%
  \tagmcbegin{artifact}%
    \hrule height 1pt%
  \tagmcend%
  \vskip 1em%
}
\makeatother
\title{Algebra First-Year Exam, May 2016, Part 1}
\author{}
\date{}
\begin{document}
\maketitle
\thispagestyle{empty}
\begin{examinstructions}
Answer four problems. (If you turn in more, the first four will be graded.) Within reason, you may quote theorems as long as you state them clearly.
\end{examinstructions}
\begin{examproblems}{0}
\item
Give an example of each of the following, with some explanation.
\begin{examsubparts}
\item[\textnormal{(a)}]
A non-solvable finite group.
\item[\textnormal{(b)}]
A non-cyclic infinite abelian group.
\item[\textnormal{(c)}]
A non-abelian solvable finite group.
\end{examsubparts}
\item
Let~\(G\) be a finite group and suppose~\(G\) acts transitively on a set~\(\Omega\). Prove that \(|\Omega|\) (the number of elements in~\(\Omega\)) is finite and that \(|\Omega|\) divides \(|G|\).
\item
Let~\(G\) be a group, and let~\(A\) and~\(B\) be normal subgroups of~\(G\) such that \(A \cap B = 1\). Prove that every element of~\(A\) commutes with every element of~\(B\).
\item
Let~\(G\) be a group of order \(14\) and let~\(H\) be a (multiplicative) cyclic group of order~\(4\). Suppose that \(\phi:G\to H\) is a group homomorphism. Let \(g\in G\). Prove that there exists some \(h\in H\) such that \(\phi(g)=h^2\).
\item
Let \(a=(1\ 2\ 3\ 4)(5\ 6)(7\ 8)\) and let \(b=(7\ 2\ 8\ 4)(1\ 6)(3\ 5)\) be elements of the symmetric group \(S_8\). Is there an element \(\sigma\in S_8\) such that \(\sigma a\sigma^{-1}=b\)? If so find one.
\item
Prove that if~\(G\) is a group of order \(56\), then~\(G\) has a normal Sylow~\(2\)-subgroup or a normal Sylow~\(7\)-subgroup.
\end{examproblems}
\end{document}
